\documentclass[justified, notoc]{tufte-handout}
\usepackage[utf8]{inputenc}
\usepackage[T1]{fontenc}
\usepackage{amsmath, amssymb}
\usepackage{graphicx}
\usepackage{minted}                     % Para resaltado de código
\usepackage{xcolor}
\usemintedstyle{colorful}
\setminted[julia]{fontsize=\footnotesize, breaklines, frame=single, framesep=2mm}

\title{Optimización del cálculo del factor de estructura mediante FFT}
\author{Usuario de Dinámica Molecular}
\date{\today}

\begin{document}

\maketitle

\begin{abstract}
El cálculo directo del factor de estructura $S(\mathbf{q})$ mediante doble suma sobre partículas se vuelve prohibitivo para sistemas grandes. La transformada rápida de Fourier (FFT) ofrece una alternativa escalable que reduce la complejidad de $\mathcal{O}(N^2)$ a $\mathcal{O}(N_g \log N_g)$. Este documento describe el fundamento matemático y proporciona una implementación práctica en Julia.
\end{abstract}

\section{Fundamento matemático}\label{sec:teoria}

El factor de estructura estático se define como
\begin{equation}
S(\mathbf{q}) = \frac{1}{N} \left\langle \left| \sum_{j=1}^{N} e^{-i\mathbf{q}\cdot\mathbf{r}_j} \right|^2 \right\rangle,
\end{equation}
donde $\mathbf{r}_j$ son las posiciones de las $N$ partículas y $\langle \cdot \rangle$ denota promedio sobre configuraciones. 
\marginnote{El vector $\mathbf{q}$ pertenece a la red recíproca de la caja de simulación: $\mathbf{q}=2\pi(n_x/L_x, n_y/L_y, n_z/L_z)$.}

La expresión anterior puede reescribirse en términos de la densidad numérica local $\rho(\mathbf{r}) = \sum_{j=1}^N \delta(\mathbf{r} - \mathbf{r}_j)$.
\marginnote{La delta de Dirac $\delta$ representa partículas puntuales.} 
Su transformada de Fourier es
\begin{equation}
\hat{\rho}(\mathbf{q}) = \int_V \rho(\mathbf{r}) e^{-i\mathbf{q}\cdot\mathbf{r}} d\mathbf{r} = \sum_{j=1}^N e^{-i\mathbf{q}\cdot\mathbf{r}_j}.
\end{equation}
Por lo tanto,
\begin{equation}
S(\mathbf{q}) = \frac{1}{N} \left\langle |\hat{\rho}(\mathbf{q})|^2 \right\rangle.
\end{equation}

Para evaluar $\hat{\rho}(\mathbf{q})$ de manera eficiente, se discretiza el espacio en una rejilla tridimensional de $M_x \times M_y \times M_z$ celdas. La densidad en cada punto de rejilla $\mathbf{r}_\mathbf{n}$ se obtiene mediante un esquema de asignación (interpolación) de las partículas. La FFT tridimensional de esta rejilla proporciona $\hat{\rho}$ en todos los vectores $\mathbf{q}$ compatibles simultáneamente.

\section{Ventajas y limitaciones}\label{sec:ventajas}

El método FFT presenta los siguientes puntos a considerar:

\begin{itemize}
    \item \textbf{Velocidad}: La FFT escala como $\mathcal{O}(M\log M)$, con $M = M_x M_y M_z$, independiente de $N$. Para sistemas con $N \gtrsim 10^4$ es notablemente más rápido que el cálculo directo.
    \item \textbf{Malla completa de $\mathbf{q}$}: Se obtienen todos los vectores de la red recíproca hasta la frecuencia de Nyquist de la rejilla, lo que facilita el promediado isotrópico.
    \item \textbf{Error de interpolación}: La discretización introduce un error que depende del esquema de asignación y de la resolución de la rejilla. Esquemas de orden superior (CIC, TSC) mitigan este efecto.
    \marginnote{El error también puede corregirse dividiendo por el factor de forma de la función de asignación.}
\end{itemize}

\section{Implementación en Julia}\label{sec:implementacion}

A continuación se presenta una implementación completa utilizando el paquete \texttt{FFTW.jl}. El código incluye un esquema de asignación \textit{Cloud-in-Cell} (CIC) de primer orden, que ofrece un buen compromiso entre precisión y simplicidad.

\subsection*{Código fuente}

\begin{minted}{julia}
using FFTW, Statistics, LinearAlgebra

"""
    structure_factor_fft(positions, box_size, grid_dims)

Calcula S(q) para una configuración usando FFT con asignación CIC.
- `positions`: matriz de tamaño (3, N) con coordenadas (x,y,z) en [0, L).
- `box_size`: tupla (Lx, Ly, Lz).
- `grid_dims`: tupla (Nx, Ny, Nz) puntos de rejilla.

Retorna:
- `q_vals`: vector de magnitudes q (ordenado).
- `Sq`: vector de S(q) promediado isotrópicamente.
"""
function structure_factor_fft(positions, box_size, grid_dims)
    Lx, Ly, Lz = box_size
    Nx, Ny, Nz = grid_dims
    N_part = size(positions, 2)
    
    # 1. Inicializar rejilla de densidad (Float64)
    ρ = zeros(Float64, Nx, Ny, Nz)
    
    # Factores de conversión: índice = floor(pos / L * N) + 1
    dx, dy, dz = Lx / Nx, Ly / Ny, Lz / Nz
    inv_dx, inv_dy, inv_dz = 1.0 / dx, 1.0 / dy, 1.0 / dz
    
    # 2. Asignación Cloud-in-Cell (CIC) a la rejilla
    for j in 1:N_part
        x, y, z = positions[1, j], positions[2, j], positions[3, j]
        
        # Coordenadas continuas en unidades de celda (0-indexadas)
        fx = x * inv_dx
        fy = y * inv_dy
        fz = z * inv_dz
        
        # Índices de la celda base
        i0 = floor(Int, fx)
        j0 = floor(Int, fy)
        k0 = floor(Int, fz)
        
        # Distancias a la esquina inferior
        wx = fx - i0
        wy = fy - j0
        wz = fz - k0
        
        # Pesos para las 8 esquinas
        w000 = (1-wx)*(1-wy)*(1-wz)
        w100 = wx    *(1-wy)*(1-wz)
        w010 = (1-wx)*wy    *(1-wz)
        w001 = (1-wx)*(1-wy)*wz
        w110 = wx    *wy    *(1-wz)
        w101 = wx    *(1-wy)*wz
        w011 = (1-wx)*wy    *wz
        w111 = wx    *wy    *wz
        
        # Convertir a índices 1-based con condiciones periódicas
        i0 = mod1(i0, Nx); i1 = mod1(i0+1, Nx)
        j0 = mod1(j0, Ny); j1 = mod1(j0+1, Ny)
        k0 = mod1(k0, Nz); k1 = mod1(k0+1, Nz)
        
        # Acumular densidad
        ρ[i0, j0, k0] += w000
        ρ[i1, j0, k0] += w100
        ρ[i0, j1, k0] += w010
        ρ[i0, j0, k1] += w001
        ρ[i1, j1, k0] += w110
        ρ[i1, j0, k1] += w101
        ρ[i0, j1, k1] += w011
        ρ[i1, j1, k1] += w111
    end
    
    # 3. Transformada rápida de Fourier
    ρ_hat = fft(ρ)
    
    # 4. Calcular |ρ_hat|^2 normalizado
    S_all = abs2.(ρ_hat) ./ N_part
    
    # 5. Construir vectores q y agrupar isotrópicamente
    #    Primero obtenemos las coordenadas de cada punto en el espacio recíproco
    qx = fftfreq(Nx, 2π/dx) |> fftshift
    qy = fftfreq(Ny, 2π/dy) |> fftshift
    qz = fftfreq(Nz, 2π/dz) |> fftshift
    
    S_shifted = fftshift(S_all)   # Centrar el origen de frecuencias
    
    q_vals = Float64[]
    S_q_vals = Float64[]
    
    # Diccionario para acumular por magnitud de q
    groups = Dict{Float64, Vector{Float64}}()
    
    for k in 1:Nz, j in 1:Ny, i in 1:Nx
        q = sqrt(qx[i]^2 + qy[j]^2 + qz[k]^2)
        if q == 0.0
            continue  # excluir el pico central
        end
        # Redondear a 6 decimales para agrupar magnitudes similares
        key = round(q, digits=6)
        push!(get!(groups, key, Float64[]), S_shifted[i,j,k])
    end
    
    for (key, vals) in groups
        q = parse(Float64, key)  # key es string por redondeo, mejor usar Float64 directamente
        # Alternativa: convertir a Float64
        q_val = key
        push!(q_vals, q_val)
        push!(S_q_vals, mean(vals))
    end
    
    # Ordenar por q
    idx = sortperm(q_vals)
    return q_vals[idx], S_q_vals[idx]
end
\end{minted}

\subsection*{Notas sobre el código}
\marginnote[-2\baselineskip]{La función \texttt{fftfreq} de FFTW genera las frecuencias sin desplazar. \texttt{fftshift} centra el origen en el centro del arreglo para facilitar la interpretación.}
\begin{itemize}
    \item El esquema CIC distribuye la contribución de cada partícula entre los 8 vértices de la celda que la contiene, proporcionalmente a la distancia.
    \item La normalización por $N$ produce $S(q)$ que tiende a $1$ para $q \to \infty$ (asumiendo ruido de disparo corregido).
    \item El promediado isotrópico agrupa todos los puntos de la rejilla recíproca con igual magnitud $q$. El número de puntos por esfera crece con $q$, lo que mejora la estadística a $q$ altos.
\end{itemize}

\section{Corrección por factor de forma de la asignación}\label{sec:correccion}

La interpolación de partículas puntuales a una rejilla actúa como una convolución con la función de asignación $W(\mathbf{r})$. En el espacio recíproco, esto multiplica el verdadero $\hat{\rho}(\mathbf{q})$ por la transformada de Fourier de $W$, es decir, $\hat{W}(\mathbf{q})$. Para recuperar $S(\mathbf{q})$ sin sesgo, se debe dividir $|\hat{\rho}_{\text{grid}}(\mathbf{q})|^2$ por $|\hat{W}(\mathbf{q})|^2$.
\marginnote{Para la asignación CIC en una rejilla ortorrómbica, el factor de corrección es $\hat{W}(\mathbf{q}) = \operatorname{sinc}\left(\frac{q_x \Delta x}{2}\right) \operatorname{sinc}\left(\frac{q_y \Delta y}{2}\right) \operatorname{sinc}\left(\frac{q_z \Delta z}{2}\right)$.}
En el código anterior no se incluye esta corrección por brevedad; en implementaciones precisas se debe aplicar antes del promediado isotrópico.

\end{document}
